\unit{A practical mathematical toolkit}
\unitref{Common methods across all four course parts}
\subsection{Expectation values in an eigenbasis}
Suppose \(\ket\psi=\sum_nc_n\ket n\),
where \(\ket n\) are eigenstates of an
observable \(A\) with eigenvalues \(a_n\).
Then
\[
\langle A\rangle=\sum_n|c_n|^2a_n,\quad
(\Delta A)^2=\sum_n|c_n|^2a_n^2
-\left(\sum_n|c_n|^2a_n\right)^2.
\]
The \emph{probabilities} are \(|c_n|^2\),
not \(c_n\). For operators not diagonal
in this basis, retain cross terms
\(c_m^*c_n\bra mA\ket n\).

\subsection{Commutators and transformations}
For Heisenberg evolution
\(\dot O=(i/\hbar)[H,O]\).
If \(H\) is a two-level Hamiltonian
\(H=(\hbar\Omega/2)\sigma_x\),
\[
U(t)=e^{-i\Omega t\sigma_x/2}
=I\cos(\Omega t/2)
-i\sigma_x\sin(\Omega t/2),
\]
because \(\sigma_x^2=I\).
For field displacement
\(D=e^{\alpha a^\dagger-\alpha^*a}\),
the identity \(D^\dagger aD=a+\alpha\)
follows by truncating the nested
commutator expansion after one step.

\subsection{Useful sums and integrals}
The geometric series for \(|q|<1\) is
\(\sum_{n=0}^\infty q^n=(1-q)^{-1}\).
Differentiate with respect to \(q\) to obtain
\[
\sum_{n=0}^\infty nq^n=\frac{q}{(1-q)^2},\quad
\sum_{n=0}^\infty n^2q^n
=\frac{q(1+q)}{(1-q)^3}.
\]
For \(\beta>0\) and integer \(m\ge0\),
\[
\boxed{\int_0^\infty r^me^{-\beta r}dr
=\frac{m!}{\beta^{m+1}}.}
\]
These formulas appear in thermal
photon statistics and hydrogen radial
expectation calculations.

\subsection{Perturbation theory}
For \(H=H_0+\lambda V\) and a
nondegenerate unperturbed state \(\ket i\),
\[
E_i^{(1)}=\bra iV\ket i,\qquad
E_i^{(2)}=
\sum_{n\ne i}\frac{|\bra nV\ket i|^2}
{E_i^{(0)}-E_n^{(0)}}.
\]
A first-order correction to the
\emph{state} is
\[
\ket{i^{(1)}}=
\sum_{n\ne i}\ket n
\frac{\bra nV\ket i}{E_i^{(0)}-E_n^{(0)}}.
\]
Use the energy formula for a
finite-nucleus or Stark \emph{shift};
use the state-mixing formula to
explain how weak parity violation
creates a forbidden transition.
For degenerate levels, first
diagonalise the perturbation in
the degenerate subspace.

\subsection{Fourier transforms of differential equations}
With the convention
\(\tilde f(\omega)=\int f(t)e^{-i\omega t}dt\),
\(\widetilde{\dot f}=i\omega\tilde f\).
Hence
\[
\dot a=-\lambda a+f(t)
\quad\Longrightarrow\quad
\boxed{\tilde a(\omega)=
\frac{\tilde f(\omega)}{\lambda+i\omega}}.
\]
A cavity susceptibility is an
instance of this elementary
linear-response formula.
The real damping rate appears
in the denominator together
with frequency detuning.

\subsection{Normalisation of conditional quantum states}
A detected branch is generally
\(\tilde\rho=K\rho K^\dagger\).
Keep its trace as the success
probability \(p=\operatorname{Tr}\tilde\rho\)
and only then write \(\rho_{\rm cond}=\tilde\rho/p\).
It is mathematically invalid to
normalise individual terms before
adding coherent amplitudes or
to sum conditional probabilities
without their branch weights.

\begin{worked}{A quick selection-rule test}
Suppose a putative E2 transition goes
from \(J=0\) to \(J'=1\).
E2 has tensor rank \(k=2\).
The triangle rule requires
\[
|0-1|\le2\le0+1,
\]
which fails at the right-hand inequality.
Therefore the rank-2 matrix element
vanishes before any radial integral
is evaluated. A parity check is
still needed for other pairs
that pass the triangle rule.
\end{worked}

\begin{checkpoint}
Often the difficult part is recognising
the \emph{right representation}.
Fock expansions for photon counts;
quadratures for homodyne/squeezing;
Kraus operators for channels;
small angular-momentum blocks for
atomic splittings; reduced matrix
elements for spectral transitions.
\end{checkpoint}
